% Draft status: V0
% Evidence status: checked where possible; TODOs mark missing evidence.
\section{Related Work}

\noindent\textbf{Visual representation paradigms and modern recognition backbones.}
The models in this paper represent four broad approaches to visual representation learning. ResNet50 is the classical residual convolutional baseline, using skip connections to train deep convolutional networks effectively \citep{he2016resnet}. ConvNeXtV2 modernizes the convolutional family through architectural and representation-learning changes, including a fully convolutional masked-autoencoder design and Global Response Normalization \citep{woo2023convnextv2}. DeiT3 revisits supervised Vision Transformer training and provides a strong ViT recipe \citep{touvron2022deit3}. SwinV2 develops a scalable hierarchical transformer with shifted-window attention and training stabilizations for large-scale vision models \citep{liu2022swinv2}. MaxViT combines convolutional structure with local-global multi-axis attention \citep{tu2022maxvit}; this paper places it under the Transformer Vision Paradigm because its central contribution is attention-based aggregation. DINOv2 and MAE provide self-supervised representation-learning baselines, with DINOv2 emphasizing large-scale self-supervised visual features \citep{oquab2024dinov2} and MAE learning through masked patch reconstruction \citep{he2022mae}. CLIP and SigLIP represent vision-language foundation models: CLIP learns transferable visual models through contrastive image-text pretraining \citep{radford2021clip}, while SigLIP modifies image-text learning with a sigmoid loss \citep{zhai2023siglip}. This paper uses these citations to describe model families, not to claim that any cited architecture paper establishes regional robustness.

\noindent\textbf{Dataset bias, domain shift, and geographic robustness.}
Dataset bias refers here to systematic skew in collection, representation, labels, or visual context. Early work on dataset bias showed that recognition datasets can contain dataset-specific cues that affect cross-dataset generalization \citep{torralba2011datasetbias}. ImageNet and its successors have also motivated questions about whether conventional benchmark progress transfers to distribution-shifted settings \citep{deng2009imagenet,recht2019imagenet,taori2020naturalshifts}. Domain-shift benchmarks and evaluation suites such as WILDS and DomainBed emphasize that models can behave differently when training and evaluation distributions differ \citep{koh2021wilds,gulrajani2021domainbed}. Geographic robustness is related but more specific: it asks whether performance remains consistent across regional groups. GeoDE directly frames object recognition as a geographically diverse evaluation problem \citep{ramaswamy2023geode}, while Dollar Street provides household-object images with geographic and socioeconomic metadata \citep{rojas2022dollarstreet}. Prior geodiversity and dataset-audit work supports the premise that geographic representation matters in visual datasets \citep{shankar2017geodiversity,wang2022revise,kalluri2023geonet}. These literatures motivate the evaluation design, but they do not by themselves explain the mechanisms behind the empirical gaps observed in this paper.

\noindent\textbf{Metrics for regional disparity and robustness evaluation.}
Robustness evaluation requires more than average accuracy. Worst-group and hidden-stratification work argues that aggregate performance can obscure weak performance for particular groups or subclasses \citep{sagawa2020groupdro,sohoni2020hiddenstratification}. Fairness-oriented vision benchmarks similarly motivate reporting performance across groups rather than treating a single aggregate as sufficient \citep{gustafson2023facet}. This paper adapts that intuition to geographic regions by reporting overall accuracy, per-region accuracy, worst-region accuracy, best-region accuracy, and the best-minus-worst geographic-disparity gap. It also reports pretrained-versus-fine-tuned deltas in both accuracy and disparity. The metric is intentionally simple: it identifies uneven regional performance, but it does not prove discrimination, causal mechanisms, or deployment harm. TODO: add confidence intervals and per-region sample counts before the final version.
